\section{Strategy Selection and User Workflow}
The estimator can answer one fixed-configuration question or search over a
small set of feasible strategies. A candidate is an \texttt{EstimatorConfig}
paired with a resolved \texttt{ModelSpec}. The CLI search path evaluates the
requested resource and distributed-mode options, separates feasible and
infeasible cases, and ranks candidates by feasibility, distributed-mode
simplicity, and memory headroom.

The simplicity order is:
\[
\texttt{tp\_only}
\prec
\texttt{ddp}
\prec
\texttt{zero2}
\prec
\texttt{zero3}.
\]
This favors a smaller distributed stack when the memory estimate permits it.
The UI uses a richer candidate set over distributed mode,
checkpointing, tensor-parallel degree, and sequence-parallel settings. DDP and
ZeRO require enough GPUs for data parallelism after tensor parallelism is
assigned, while sequence parallelism requires tensor parallelism.

Candidate generation follows the same constraints a user would apply by hand.
The resource pool determines the available GPU count and memory budget. Tensor
parallelism consumes GPUs inside each model replica. Data parallelism then uses
the remaining replica count. ZeRO stages are considered only when a data
parallel group exists. Sequence parallelism is paired with tensor parallelism,
since it partitions sequence-local activation state across the tensor-parallel
group.

\begin{center}
\small
\begin{tabular}{@{}p{0.25\linewidth}p{0.65\linewidth}@{}}
\toprule
Strategy family & Search treatment \\
\midrule
Single GPU & Used when only one GPU is specified. \\
Tensor parallel & Shards eligible matrices and associated tensor-local state
across the tensor-parallel group. \\
DDP & Replicates model state and changes communication workspace through
reducer buckets. \\
ZeRO-2 & Shards gradients and optimizer state across data-parallel ranks. \\
ZeRO-3 & Shards trainable parameters as well, adding gather and prefetch
windows around computation. \\
LoRA & Changes trainable tensors and optimizer state through adapter rank and
target modules. \\
\bottomrule
\end{tabular}
\normalsize
\end{center}

Among feasible results, the reported headroom gives the user a practical margin against
allocator variation, dataloader details, and small differences between the
synthetic estimate and the real training script. The search output therefore
contains both the selected candidate and the alternatives, allowing a user to
choose a slower option with more margin when the cluster allocation is scarce.

\subsection{Slowdown Model}
Slowdown estimation is separate from memory estimation. The web recommendation
path uses a small heuristic prior:
\[
R_{\mathrm{slow}} =
\left(1 + s_{\mathrm{dist}}\right)
\left(1 + s_{\mathrm{ckpt}}\right)
\left(1 + s_{\mathrm{tp}}\right)
\left(1 + s_{\mathrm{sp}}\right) - 1.
\]
The distributed-mode priors are $0\%$ for single-GPU and pure tensor
parallelism, $3\%$ for DDP, $8\%$ for ZeRO-2, and $18\%$ for ZeRO-3.
Checkpointing contributes $30\%$ because recomputation increases backward
work. Tensor parallelism contributes $0.08\log_2(T_{\mathrm{P}})$, and
sequence parallelism contributes $4\%$ when enabled.

The slowdown score is simple because estimating training time is a hard problem
outside the scope of this project. It expresses the expected slowdown
relative to the least-sharded feasible candidate.
For example, checkpointing receives a large penalty because it trades memory for
extra backward compute. ZeRO-3 receives a larger distributed penalty than
ZeRO-2 because parameter gathering introduces additional synchronization and
workspace.

\subsection{Training Export}
Search results can be exported into a TRL-compatible YAML strategy. The export
contains the model reference, tuning mode, optimizer and precision choices,
sequence length, checkpointing, attention backend, distributed strategy, and
LoRA fields when applicable. The training runner consumes the first exported
candidate by default. In the intended workflow, the user searches over model,
context length, resource pool, and distributed strategy, exports the best
candidate, and launches training from that file:

\begin{verbatim}
simplesft search <HF Model> ... --export-file ./trl_strategy.yaml
simplesft train --config ./trl_strategy.yaml --dataset <HF Dataset>
\end{verbatim}

The command-line interface and web server return the same core objects: peak
memory, feasibility, phase records, component breakdowns, assumptions, and
diagnostic metadata.
